% Part: first-order-logic
% Chapter: syntax-and-semantics
% Section: substitution

\documentclass[../../../include/open-logic-section]{subfiles}

\begin{document}

\olfileid[tr]{fol}{syn}{ext}

\olsection{Kaplamsallık}

\begin{explain}
Bazen ilgililik de denen kaplamsallık, gayriresmî olarak şöyle ifade
edilebilir: !!{formula}~$!A$ ifadesinin !!a{structure}~$\Struct M$
içinde !!a{variable} ataması~$s$'ye göre sağlanmasını etkileyen
tek etmenler, !!{domain} büyüklüğü ile $\Struct M$ ve $s$ tarafından
$!A$ içinde gerçekten geçen dil ögelerine yapılan atamalardır.

Kaplamsallığın doğrudan bir sonucu şudur: İki !!{structure}s~$\Struct
M$ ve $\Struct M'$ aynı evrene sahipse ve $!A$ tümcesinde geçen bütün
dil ögeleri üzerinde uyuşuyorsa $!A$'nın doğru olup olmadığı konusunda
da uyuşmak zorundadır.
\end{explain}

\begin{prop}[Kaplamsallık]
  \ollabel{prop:extensionality}
  $!A$ !!a{formula}; $\Struct M_1$ ve $\Struct M_2$,
  $\Domain{M_1}=\Domain{M_2}$ olan !!{structure}s; $s$ ise ortak
  evren üzerinde bir değişken ataması olsun. $!A$ içinde geçen her
  !!{constant}~$c$, bağıntı simgesi~$R$ ve !!{function}~$f$ için
  $\Assign{c}{M_1}=\Assign{c}{M_2}$,
  $\Assign{R}{M_1}=\Assign{R}{M_2}$ ve
  $\Assign{f}{M_1}=\Assign{f}{M_2}$ ise
  $\Sat{M_1}{!A}[s]$ ancak ve ancak $\Sat{M_2}{!A}[s]$.
\end{prop}

\begin{proof}
  Önce her terim için $\Value{t}{M_1}[s]=\Value{t}{M_2}[s]$
  eşitliğini $t$ üzerinde tümevarımla ispatlayın. Ardından, tümevarımın
  $!A$'nın atomik olduğu temel durumunda az önce ispatlanan savı
  kullanarak önermeyi $!A$ üzerinde tümevarımla ispatlayın.
\end{proof}

\begin{prob}
\olref[fol][syn][ext]{prop:extensionality} önermesinin ispatını
ayrıntılı olarak verin.
\end{prob}

\begin{cor}[!!^{sentence}s için kaplamsallık]
  \ollabel{cor:extensionality-sent}
  $!A$ !!a{sentence}; $\Struct M_1$ ve $\Struct M_2$ ise
  \olref{prop:extensionality} önermesindeki gibi olsun. O zaman
  $\Sat{M_1}{!A}$ ancak ve ancak $\Sat{M_2}{!A}$.
\end{cor}

\begin{proof}
Sonuç, \olref{prop:extensionality} önermesinden
\olref[ass]{cor:sat-sentence} aracılığıyla çıkar.
\end{proof}

Ayrıca bir terimin değeri ve !!a{structure} ifadesinin !!a{formula}
ifadesini sağlayıp sağlamadığı yalnızca alt terimlerin değerlerine
bağlıdır.

\begin{prop}\ollabel{prop:ext-terms}
$\Struct M$ !!a{structure}; $t$ ile $t'$ terimler; $s$ bir değişken
ataması olsun. O zaman
$\Value{\Subst{t}{t'}{x}}{M}[s]=
\Value{t}{M}[\Subst{s}{\Value{t'}{M}[s]}{x}]$.
\end{prop}

\begin{proof}
$t$ üzerinde tümevarımla.
\begin{enumerate}
\item $t$ bir sabitse, örneğin $t\ident c$ ise
  $\Subst{t}{t'}{x}=c$ ve
  $\Value{c}{M}[s]=\Assign{c}{M}=
  \Value{c}{M}[\Subst{s}{\Value{t'}{M}[s]}{x}]$.

\item $t$, $x$ dışında bir değişkense, örneğin $t\ident y$ ise
  $\Subst{t}{t'}{x}=y$ ve
  $\varAssign{s}{\Subst{s}{\Value{t'}{M}[s]}{x}}{x}$ olduğundan
  $\Value{y}{M}[s]=
  \Value{y}{M}[\Subst{s}{\Value{t'}{M}[s]}{x}]$.

\item $t\ident x$ ise $\Subst{t}{t'}{x}=t'$ olur. Öte yandan
  $\Subst{s}{\Value{t'}{M}[s]}{x}$ tanımı gereği
  $\Value{x}{M}[\Subst{s}{\Value{t'}{M}[s]}{x}]
  =\Value{t'}{M}[s]$.

\item $t\ident\Atom{f}{t_1,\dots,t_n}$ ise:
\begin{multline*}
  \Value{\Subst{t}{t'}{x}}{M}[s]  = \\
  \begin{aligned}[b]
& = \Value{\Atom{f}{\Subst{t_1}{t'}{x}, \dots, \Subst{t_n}{t'}{x}}}{M}[s]\\
& \qquad    \text{$\Subst{t}{t'}{x}$ tanımı gereği}\\
& = \Assign{f}{M}(\Value{\Subst{t_1}{t'}{x}}{M}[s], \dots,
    \Value{\Subst{t_n}{t'}{x}}{M}[s])\\
& \qquad  \text{$\Value{\Atom{f}{\dots}}{M}[s]$ tanımı gereği}\\
& = \Assign{f}{M}(\Value{t_1}{M}[\Subst{s}{\Value{t'}{M}[s]}{x}], \dots,
   \Value{t_n}{M}[\Subst{s}{\Value{t'}{M}[s]}{x}])\\
& \qquad    \text{tümevarım varsayımı gereği}\\
& = \Value{t}{M}[\Subst{s}{\Value{t'}{M}[s]}{x}]
    \text{, $\Value{\Atom{f}{\dots}}{M}[\Subst{s}{\Value{t'}{M}[s]}{x}]$ tanımı gereği.}
  \end{aligned}
\end{multline*}
\end{enumerate}
\end{proof}

\begin{prop}\ollabel{prop:ext-formulas} $\Struct M$ !!a{structure},
$!A$ !!a{formula}, $t'$ bir terim ve $s$ bir değişken ataması olsun.
O zaman $\Sat{M}{\Subst{!A}{t'}{x}}[s]$ ancak ve ancak
$\Sat{M}{!A}[\Subst{s}{\Value{t'}{M}[s]}{x}]$.
\end{prop}

\begin{proof}
Alıştırma.
\end{proof}

\begin{prob}
\olref[fol][syn][ext]{prop:ext-formulas} önermesini ispatlayın.
\end{prob}

\begin{explain}
  \cref{fol:syn:ext:prop:ext-terms,fol:syn:ext:prop:ext-formulas}
  sonuçlarının önemi şudur. Bir $t$ terimi ya da !!a{formula}~$!A$ ile
  bir $t'$ terimimiz olsun; $\Subst{t}{t'}{x}$ değerini ya da
  $\Subst{!A}{t'}{x}$ ifadesinin !!a{structure}~$\Struct M$ içinde
  !!a{variable} ataması~$s$'ye göre sağlanıp sağlanmadığını öğrenmek
  isteyelim. Önce yerine koymayı yapıp ardından $\Struct M$ ve $s$'ye
  göre değeri veya sağlanmayı inceleyebiliriz. Ya da önce $t'$ teriminin
  $\Struct M$ içinde $s$'ye göre $m=\Value{t'}{M}[s]$ değerini bulup
  !!{variable} atamasını $\Subst{s}{m}{x}$ olarak değiştirebilir ve
  sonra $t$'nin $\Struct M$ ile $\Subst{s}{m}{x}$ içindeki değerine ya
  da $\Sat{M}{!A}[\Subst{s}{m}{x}]$ olup olmadığına bakabiliriz.
  \Cref{fol:syn:ext:prop:ext-terms,fol:syn:ext:prop:ext-formulas},
  hangi yolu seçersek seçelim yanıtın aynı olacağını güvence altına alır.
\end{explain}

\end{document}
